% Coverage certificate for heart-s5-c2 (genome 57c311c6-56f5-4c5d-a475-46433ef58bf2); content sha256 e5293773fce87faec815a166111f195edf55b68bde14ba046b51556ab442dcea.
\providecommand{\explaneatcertdir}[1]{}
\begin{drafted}[heart-s5-c2 coverage certificate]
\section{Coverage certificate: heart-s5-c2}\label{app:heart-s5-c2:sec}
Provenance, coverage and checks of the explanation, rendered from the canonical JSON record \texttt{heart-s5-c2.json} (schema 1.4.0, SHA-256 \texttt{e5293773fce87fae}\ldots). The explanation itself (summary, parts, evidence, figures) is the read-down document before this page.

\begin{table}[htbp]
\centering\small
\begin{tabularx}{\linewidth}{@{}lX@{}}
\toprule
Dataset & Heart Disease (binary, prepared) \\
Experiment & 10seed-heart-s5, seed 5 \\
Genome & \texttt{57c311c6-56f5-4c5d-a475-46433ef58bf2} (NEAT key 9858, generation 39) \\
Explanation & \texttt{dabdb4bb-e227-46eb-a995-d4123373f234} (39 operations) \\
Train AUC / accuracy & 0.9389 / 0.8802 (n = 242) \\
Test AUC / accuracy & 0.8922 / 0.8197 (n = 61) \\
Output activation & sigmoid \\
Phenotype & 32 nodes (26 inputs, 5 hidden, 1 output), 35 enabled edges \\
Explained model state & 41 nodes, 37 enabled edges \\
Refactoring operations & add\_\allowbreak{}identity\_\allowbreak{}node 2, split\_\allowbreak{}node 5 \\
Function preservation & max $|\Delta\,\mathrm{output}|$ = $2.22\times 10^{-16}$ on the training split \\
Annotations & 6 (5 leaf, 1 composition) \\
Structural coverage & 1 (40 of 40 candidate nodes; 5 Null) \\
Compositional coverage & 1 \\
Uncovered nodes / edges & 0 / 0 \\
Missing compositions & 0 \\
Full coverage & yes \\
Drafted by & AI-assisted draft \\
Verified by & unverified \\
Record hash (SHA-256) & \texttt{e5293773fce87fae}\ldots \\
\bottomrule
\end{tabularx}
\caption{Explanation certificate summary: heart-s5-c2.}
\label{app:heart-s5-c2:summary}
\end{table}
\subsection{Coverage}\label{app:heart-s5-c2:coverage}
Coverage is computed by the current code (the function behind the \texttt{get\_coverage} tool) under these conventions. \emph{Candidate nodes} are all nodes of the explained model state except the output nodes. A candidate node is \emph{covered} when some annotation of the hierarchy (a leaf, or a composition together with its descendants) contains it and all its enabled outgoing edges. \emph{Exact Null} parts, which have no effect on any output for any input, count as covered: nodes with no path to an output over enabled, non-zero-weight edges (disconnected inputs and dead ends), and edges that have zero weight or touch such a node. Compositional coverage is the fraction of required compositions present: every annotation except a single global root containing every output must have a parent; a parentless, childless annotation made only of Null nodes is exempt.

\begin{itemize}
\item \textbf{Null, disconnected input (5):} \texttt{-13} (restecg:1), \texttt{-20} (slope:2), \texttt{-26} (thal:7), \texttt{-3} (sex:1), \texttt{-5} (cp:2).
\item \textbf{Uncovered nodes:} none.
\item \textbf{Uncovered edges:} none.
\item \textbf{Missing compositions:} none.
\end{itemize}
\noindent\textbf{Agreement with the \texttt{get\_coverage} tool:} yes.

\subsection{Closed-form checks}\label{app:heart-s5-c2:closedform}
{\small\noindent\begin{tabularx}{\linewidth}{@{}lXrr@{}}
\toprule Tag & Annotation & Closed form vs forward pass & Exact function vs forward pass \\ \midrule
A1 & heart\_\allowbreak{}disease\_\allowbreak{}logit & 0.00243 & $2.22\times 10^{-16}$ \\
A2 & ischaemia\_\allowbreak{}gate & 0.0173 & 0 \\
A3 & bp\_\allowbreak{}offset\_\allowbreak{}gate & 0.00119 & $1.78\times 10^{-15}$ \\
A4 & ischaemia\_\allowbreak{}damping\_\allowbreak{}thal6 & 0.00134 & $2.22\times 10^{-16}$ \\
A5 & lvh\_\allowbreak{}without\_\allowbreak{}exercise\_\allowbreak{}angina & $2.52\times 10^{-4}$ & 0 \\
A6 & high\_\allowbreak{}fasting\_\allowbreak{}sugar & $3.12\times 10^{-4}$ & 0 \\
\bottomrule
\end{tabularx}\par}
\subsection{Refactoring operations}\label{app:heart-s5-c2:operations}
The explained model state is the phenotype after these operations, replayed in order. All are function-preserving; the largest change in the model output over the 242 training rows is $2.22\times 10^{-16}$ (tolerance $1.00\times 10^{-9}$).
\begin{itemize}
\item \#26: split \texttt{-12} into \texttt{-12\_\allowbreak{}a}, \texttt{-12\_\allowbreak{}b}.
\item \#27: split \texttt{-17} into \texttt{-17\_\allowbreak{}a}, \texttt{-17\_\allowbreak{}b}, \texttt{-17\_\allowbreak{}c}.
\item \#28: split \texttt{-25} into \texttt{-25\_\allowbreak{}a}, \texttt{-25\_\allowbreak{}b}, \texttt{-25\_\allowbreak{}c}.
\item \#29: split \texttt{-8} into \texttt{-8\_\allowbreak{}a}, \texttt{-8\_\allowbreak{}b}.
\item \#30: split \texttt{-6} into \texttt{-6\_\allowbreak{}a}, \texttt{-6\_\allowbreak{}b}.
\item \#31: identity node \texttt{identity\_\allowbreak{}290} before \texttt{290} (1 connection(s)).
\item \#32: identity node \texttt{identity\_\allowbreak{}1416} before \texttt{1416} (1 connection(s)).
\end{itemize}
\subsection{Audit against the definitions}\label{app:heart-s5-c2:audit}
Automatic findings on the stored explanation (codes as in the JSON record's \texttt{issues}); they flag what an examiner should check, not verdicts.
\begin{itemize}
\item \textsc{warning} \texttt{unverified}: No human verification is recorded for this explanation.
\end{itemize}
\end{drafted}
